\section{Method}
In this section, we present \textsc{ESAT}, our environment-free pipeline that transforms API specifications into a dataset of multi-step, API-calling trajectories without relying on real API execution, backend infrastructure, or human-written tasks. As illustrated in \Cref{fig:pipeline}, the \textsc{ESAT} pipeline consists of three stages: (i) task synthesis from API specifications, (ii) trajectory synthesis using a teacher agent paired with an LLM-based API simulator, and (iii) trajectory filtering via an LLM judge.


\subsection{Task Synthesis}
This stage generates diverse user tasks solvable using the available APIs. The sole required input is a set of API specifications detailing each API's name, purpose, and input/output schemas.

\paragraph{Bucketized generation.} To ensure broad and balanced coverage, task generation is organized into a combinatorial grid of \emph{configuration buckets}. Each bucket is defined by a tuple of attributes, including \emph{difficulty level} (\eg, easy, medium, hard), \emph{action type} (\eg, read, write, mixed), \emph{task focus} (\eg, constraint satisfaction, derivation), the \emph{number of apps} (1, 2, or 3), and a \emph{range of APIs per app}. These buckets are processed in a randomized sequence, and a task generator LLM is prompted multiple times to yield a desired number of tasks per bucket. During each invocation, we randomly sample apps matching the bucket's required number and provide their full API specifications to the LLM, alongside the target difficulty, action type, task focus, and API count. The LLM then generates a task along with the list of APIs required to solve it. Detailed descriptions of the configuration attributes and their values are provided in \Cref{sec:bucket-gen}.

\paragraph{Inverse frequency-based sampling.} To ensure uniform coverage of apps and APIs, we employ an inverse frequency-based sampling strategy. We maintain running usage counts for every app and API across all generated tasks. When selecting apps for a generation step, we sample them with probabilities inversely proportional to their current counts, thereby prioritizing underrepresented apps. Furthermore, we rank all available APIs within the selected apps based on their counts, and explicitly instruct the task generator LLM to focus on the least-used APIs.

\paragraph{Task filtering.} Each generated task is evaluated by an LLM judge to ensure it is well-formed, solvable, and consistent with the corresponding bucket configuration.

\paragraph{Task rewriting.} Tasks generated directly from API specifications can sometimes suffer from procedural verbosity, outlining detailed instructions rather than natural goals. To address this, all tasks accepted by the judge are subsequently rewritten by an LLM to compress any procedural scaffolding into concise, intent-level requests that reflect actual user needs. These rewritten tasks are then verified by an LLM judge to ensure they remain valid. \Cref{sec:task_rewriting} provides examples illustrating the impact of this rewriting process. If a rewritten task is rejected by the judge, the corresponding initial task is retained for trajectory synthesis.

\subsection{Trajectory Synthesis}
In this stage, a teacher LLM agent iteratively solves each task by issuing API calls, while an LLM-based API simulator dynamically generates synthetic responses. This yields complete, multi-step training trajectories that include the agent's reasoning, API calls, and simulated environmental observations, all without executing any real APIs. Example trajectories are provided in \Cref{app:trajectories}.

\looseness=-1
To ensure structural correctness and semantic coherence, the simulator LLM conditions its output on a rich input context, which includes: (1) the full API specification, (2) the API call arguments, (3) the task being solved, (4) the history of previous API calls in the trajectory and their simulated responses, (5) the name, email, and phone number of a virtual user, and (6) the current date and time. To prevent overly long input contexts and accelerate simulation, we restrict the simulation history to previous API calls from the same app.


\paragraph{Rules for the simulator.} The simulator LLM is guided by four core rules encoded in its system prompt: (1) \emph{Data consistency}: it must preserve state across turns (\eg, a song assigned to a specific artist earlier must remain under that artist); (2) \emph{Data diversity}: listed results should contain a mix of task-relevant and irrelevant items to mirror real-world noise; (3) \emph{Data realism}: generated values must respect field constraints such as numeric ranges and unique identifiers; and (4) \emph{Schema compliance}: responses must strictly conform to the schema provided in the API specification.

\paragraph{Simulation quality checks.} To ensure that the simulator produces accurate and realistic responses, we employ multiple stages of quality control. Before invoking the simulator LLM, the API call arguments are programmatically validated for missing parameters and type mismatches. If errors are detected, the simulator directly returns an appropriate error message without invoking the LLM. For valid API calls, the LLM-generated response undergoes multiple checks. First, it is programmatically verified against the expected output schema. In the event of a schema mismatch, the simulator LLM is re-prompted using the schema errors as feedback. Once the response passes the schema check, it is evaluated for semantic correctness and consistency with the simulation history using the same LLM as a judge. If the judge identifies any errors, the simulator LLM is prompted to refine its response using the judge's output as feedback. Only responses that pass all checks are returned. We allow a fixed number of refinement steps; if the simulator fails to generate a valid response after exhausting all retries, we declare a simulation failure.

\paragraph{Trajectory termination.} Trajectory generation for a task concludes when the agent explicitly calls a special task completion API, exceeds a predefined limit on the number of execution steps, or invokes an API from an app that was not used to synthesize the corresponding task. The process also halts if the simulator fails.

\subsection{Trajectory Filtering}
Trajectories in which the agent finished the task by calling the task completion API are finally evaluated using a trajectory-level LLM judge, which decides whether the agent correctly solved the corresponding task. Each trajectory is judged multiple times, and we retain only those with a majority-positive verdict.

All the LLM prompts used for data synthesis are provided in~\Cref{sec:simulator-prompt}.

